\subsection{Workflow Execution}

In execution, the pipeline produces one prepared dataset per antibiotic per scope. The final prepared datasets are stored separately for the \texttt{strict}, \texttt{broad}, and \texttt{all} settings. This separation is deliberate. Data preparation is completed before model training, so the training workflow no longer performs raw TSV ingestion or hidden preprocessing. Instead, the modeling stage consumes only prepared antibiotic-specific tables. This design makes the experimental comparison across scopes reproducible, prevents hidden data-processing drift between runs, and keeps biological preprocessing separate from model-training logic. For each scope, the modeling stage writes a consistent set of outputs, namely per-antibiotic evaluation metrics, dataset summaries including feature counts and zero-feature rows, the run configuration, and ranked feature-importance scores. This execution design is important for the report because it means every result in the experiments chapter can be linked back to one prepared input scope, one antibiotic, and one consistent training configuration, and the final workflow therefore supports both reproducible experimentation and biological interpretation without requiring code-level detail inside the report itself.
